\section{Introduction}
\label{sec:intro}

The Qur’an, the holy book of Islam, occupies a central role in the spiritual, intellectual, and cultural lives of nearly two billion Muslims worldwide. Revealed in Classical Arabic (Fusha), it is not only revered as sacred scripture but also preserved through its oral tradition of recitation. Mastery of recitation is guided by tajweed which is a set of phonological and prosodic rules that ensure the Qur’an is read with precision, beauty, and correctness. For Muslims, engaging with the Qur’an involves more than simply reading the text but is an act of devotion that requires both accurate pronunciation and a deep understanding of its meaning.

Despite its origins in Arabic, the Qur’an is recited and studied across every corner of the globe, including by millions of non-Arabic speakers from diverse linguistic and cultural backgrounds. This global engagement highlights the challenges faced by learners who must navigate not only the Arabic script but also its sounds, phonology, and interpretive meanings. While Islamic scholars, including Ulama (religious scholars), Fuqaha (jurists), and researchers, have made significant contributions to Qur’anic studies, the integration of modern computational approaches, particularly artificial intelligence, remains very limited compared to the rapid advances in AI. AI has the potential to illuminate linguistic patterns, aid recitation training, and support both Arabic and non-Arabic learners in mastering tajweed and understanding the Qur’an’s message.

Existing resources, however, often fall short of capturing the Qur’an’s full multimodal nature. Some datasets provide the Arabic text alone, others include translations or transliterations, and a few extend to audio recordings at either the verse or word level. Rarely are these modalities combined in a way that allows fine-grained analysis of both the written and spoken Qur’an. This fragmentation restricts their utility for advancing AI-driven research in areas such as natural language processing, speech recognition, and text-to-speech synthesis. Table \ref{tab:comparison_final} summarizes key features of prior Qur’anic datasets, where \textsc{Qur’an-MD} offers the most comprehensive multimodal coverage.

In this work, we curate and harmonize resources from three data sources to construct \textsc{Qur’an-MD}, a comprehensive multimodal dataset of the Qur’an. Our dataset unifies Arabic text, English translation, and phonetic transliteration at both the verse and word levels, while pairing each word with an aligned audio recording of its pronunciation. We provide high-quality recitations from 30 distinct reciters at the verse level, representing diverse styles and dialectical nuances (see Table \ref{tab:dataset_stats} in the Appendix). Through careful validation and consistency checks, we ensure the reliability of the textual and audio components. By making this dataset available on Hugging Face, we aim to provide scholars, linguists, and AI researchers with a standardized resource that bridges textual, phonological, and semantic dimensions. This contribution supports the study of Qur’anic language and recitation traditions. It enables and promotes future work in cutting-edge computational applications at the intersection of speech, language, and Quranic studies.